\begin{center}
\thispagestyle{empty}
\vspace{2cm}
\LARGE{\textbf{RESUMEN}}\\[1.0cm]
\end{center}
\thispagestyle{empty}
Este es el manual técnico para el proyecto de Lector PDF Minimalista. Su propósito es familiarizar al desarrollador con el proyecto y su diseño. Este proyecto nace de la necesidad de un lector de documentos pdf que no sea intrusivo y no sature la vista. Está necesidad fue expresada por el cliente Rodolfo Mora, él es quien recibe el producto final. Para solucionar este problema el objetivo general del proyecto es desarrollar una aplicación de escritorio para Windows que permita visualizar y modificar documentos pdf mediante una interfaz limpia, minimalista y clara.

La aplicación tiene funcionalidades de modificación de pdf (eliminar, reordenar, concatenar páginas, etc), visualización de pdf, presentación de pdf, entre otras. La implementación utiliza el framework Electron junto con bibliotecas especializadas para el manejo de pdf (PDF.js, pdf-lib). La arquitectura de la aplicación es modular, se separan las funciones en módulos independientes. El manual cuenta con instrucciones para descargar el código y desarrollar, para compilar y para desplegar en el entorno final. Al final, se comparte el aprendizaje acumulado junto con las problemáticas y limitaciones de la aplicación.\\
\textbf{Palabras clave: }Lector PDF, estilo minimalista.